\section{Zadanie 7 (Hubert Jastrzębski)}
Niech $X_1, \ldots, X_n$ będą niezależnymi zmiennymi losowymi z rozkładu wykładniczego z parametrem $\lambda = 1$.
Wyznacz wartość oczekiwaną k-tej wartości w posortowanym ciągu tych zmiennych.

\subsection*{Rozwiązanie}
\textbf{Lemat:} $\big(X_{(1)} = min(X_1, \ldots, X_n) \big) \sim Exp(n \lambda)$

\textbf{Dowód:}

Wiemy, że $X_1, \ldots, X_n \sim Exp(\lambda)$ i są niezależne.
Dystrybuanta każdej z nich to: $$F_{X_i}(x) = P(X_i \leq x) = 1 - e^{-\lambda x}$$
W takim razie $P(X_i \geq x) = e^{-\lambda x}$

$$P(X_{(1)} \geq x) = P(\min(X_1, \ldots, X_n) \geq x) = P(X_1 \geq x, \ldots, X_n \geq x)$$
Zmienne są niezależne, więc:
$$P(X_{(1)} \geq x) = \prod_{i=1}^n P(X_i \geq x) = (e^{-\lambda x})^n = e^{-n \lambda x}$$
Zatem $X_{(1)} \sim Exp(n \lambda)$, ckd.

\textbf{Nasze zadanie}

Z lematu wiemy, że $X_{(1)} = \min(X_1, \ldots, X_n) \sim Exp(n)$. 
Zatem $E[X_{(1)}] = \frac{1}{n}$.

Zdefiniujmy sobie teraz odstępy między kolejnymi wartościami w posortowanym ciągu:
$$Y_1 = X_{(1)}, \quad Y_2 = X_{(2)} - X_{(1)}, \quad \ldots, \quad Y_n = X_{(n)} - X_{(n-1)}$$

Ze względu na brak pamięci rozkładu wykładniczego, 
po wystąpieniu $X_{(1)}$, pozostałe $n-1$ zmiennych losowych 
``startują od nowa'' i nadal są niezależne 
i mają rozkład wykładniczy z parametrem $\lambda = 1$.

Zatem $Y_2$ to czas oczekiwania na drugą najmniejszą wartość po wystąpieniu pierwszej i ``resetowaniu'' pozostałych zmiennych,
czyli czas oczekiwania na minimum spośród $n-1$ zmiennych losowych z rozkładu wykładniczego z parametrem $\lambda = 1$.
Z lematu więc $Y_2 \sim Exp(n-1)$ i $E[Y_2] = \frac{1}{n-1}$.

Analogicznie jest dla pozostałych odstępów:
$$Y_i \sim Exp(n-i+1) \quad \implies \quad E[Y_i] = \frac{1}{n-i+1}$$

Możemy wyrazić $X_{(k)}$ przez sumę odstępów:
$$X_{(k)} = \sum_{i=1}^k Y_i \implies E[X_{(k)}] = \sum_{i=1}^k E[Y_i] = \sum_{i=1}^k \frac{1}{n-i+1}$$
$$E[X_{(k)}] = \sum_{i=0}^{k-1} \frac{1}{n-i} = \sum_{i=n-k+1}^{n} \frac{1}{i} = \sum_{i=1}^n \frac{1}{i} - \sum_{i=1}^{n-k} \frac{1}{i} = H_n - H_{n-k}$$


\subsection*{Odpowiedź}
\[
\boxed{E[X_{(k)}] = H_n - H_{n-k}}
\]
